\silentchapter{1 Formverda er alt som er tilfelle}{formverda er alt som er tilfelle}

\prop{1}{}{Formverda er alt som er tilfelle.}

Kvifor ser ein stol ut som han gjer? Kvifor ser han ikkje ut på ein annan måte? Spørsmålet verkar enkelt; svaret har vist seg å vere noko av det mest innfløkte i designteorien. Det dominerande svaret i over hundre år har vore ein variant av Louis Sullivans formulering frå 1896: forma følgjer funksjonen. Ein stol ser ut som han gjer fordi han er laga for å sitjast i. Ein bygning ser ut som han gjer fordi han er laga for å buast i. Formelen har vore så sjølvsagd at han knapt har vorte diskutert.

Jan Michl har vist at formelen er logisk tom. Funksjonen «å sitje» avgrensar ikkje forma til stolen; ho avgrensar ikkje eingong klassen til stolar. Ein krakk, ei seng, ein stein og ein jordhug tillèt alle sitting. Funksjonen er konstant; forma varierer enormt. Det som varierer, kan ikkje forklarast av det som er konstant. Michl konkluderte: form følgjer ikkje funksjon; form følgjer form. Kvar ny gjenstand er ein modifikasjon av ein eksisterande gjenstand. Denne innsikta rydda grunnen, men ho gav ikkje noko nytt å byggje på. Me står att med spørsmålet: kvifor ser ein stol ut som han gjer?

Denne traktaten gjev eit svar. Svaret byggjer ikkje på éin faktor, men på eit samspel av krefter, agentar og historiske vilkår som saman produserer det me kallar form. Svaret let seg samanfatte i éi setning:

\begin{overgang}
Form er ein posisjon i eit rom av moglege former, forma av samtidige seleksjonstrykk, navigert av agentar på fleire skalaer.
\end{overgang}

For å forstå denne setninga treng me eit presist vokabular. Det er oppgåva til dette kapittelet.

\section*{Konfigurasjonar, ikkje ting}

\prop{1.1}{o}{Formverda er totaliteten av realiserte konfigurasjonar, ikkje ting. Alt som har form, har eksakt denne forma. At noko tek éi form framfor ei anna, krev ei forklaring.}

Det fyrste steget er å skilje forma frå tingen. Ein stol er ikkje berre «ein stol»; han er ei spesifikk samansetjing av delar i ein spesifikk geometri. Endrar du vinkelen på ryggen med tre gradar, har du ein annan konfigurasjon. Endrar du materialet frå eik til bøk, har du ein annan konfigurasjon. Kvar endring, kor lita ho enn er, produserer ein ny posisjon i det rommet av moglege former me snart skal definere.

Denne presiseringa er ikkje pedanteri. Ho er naudsynleg fordi det finst ein systematisk tendens i designteori til å snakke om «stolen» som om det var éin ting. Det finst ikkje éin stol. Det finst tusenvis av konfigurasjonar som alle tilbyr sitting til ein menneskekropp, og kvar av dei okkuperer ein unik posisjon i eit fleirdimensjonalt rom av eigenskapar. Spørsmålet er ikkje kvifor «stolen» ser ut som han gjer, men kvifor *denne bestemte konfigurasjonen* vart realisert framfor dei uendeleg mange andre som var moglege.

At noko tek éi form framfor ei anna, krev ei forklaring. Denne boka er eit forsøk på å gje den forklaringa ein struktur.

\section*{Objekt og konfigurasjon}

\prop{1.2}{d}{Eit \textbf{objekt} er den enklaste bestanddelen i ei form. Objektet er udeleleg; det utgjer substansen i formverda. Alle moglege former er gjevne med objekta.}

\prop{1.21}{d}{Ein \textbf{konfigurasjon} er ei samansetjing av objekt. Forma er strukturen til konfigurasjonen.}

Kvart formelt system treng eit nedste nivå: dei primitiva som ikkje kan dekomponerast vidare utan at systemet bryt saman. I kjemi er det grunnstoffa. I lingvistikk er det fonema. I formlæra er det objektet.

Kva som tel som eit objekt, avheng av analysenivået. For ein stolmakar er objektet kanskje eit stolbein, eit sete, ein rygg. For ein materialvitar er det eit tverrsnitt av fiberretninga i eit trestykke. For ein arkitekt er det ein vegg, eit rom, ein opning. Objektet er det nivået der vidare dekomponering øydelegg den kausale kapasiteten me er interesserte i. Under det nivået finst det framleis fysikk og kjemi, men ikkje formgjeving.

Ein konfigurasjon er ei samansetjing av objekt. Forma er strukturen til denne samansetjinga; ikkje kva delar ho inneheld, men korleis dei er ordna. To stolar kan ha same delar (fire bein, eit sete, ein rygg) men ulik form fordi delane er ordna ulikt. Forma er relasjonane, ikkje elementa.

\section*{Formrommet}

\prop{1.3}{d}{\textbf{Formrommet} (\textit{morphospace}) til ein klasse er mengda av alle hennar moglege konfigurasjonar. Kvart realisert objekt okkuperer eitt punkt i dette rommet. Det empiriske formrommet er alltid ein projeksjon; forma eksisterer uavhengig av målinga.}

Formrommet er det sentrale omgrepet i denne traktaten. Det er eit matematisk rom der kvar akse representerer ein målbar eigenskap; til dømes høgde, breidde og djupn. Kvar stol som nokon gong er laga, okkuperer eitt punkt i dette rommet. Kvar stol som kunne ha vore laga men aldri vart det, okkuperer eit anna punkt. Totaliteten av alle moglege punkt er formrommet til klassen «stolar».

Omgrepet er ikkje nytt. David Raup innførte det i paleontologien i 1966 då han viste at alle moglege skjellformer kan representerast som punkt i eit tredimensjonalt rom definert av tre geometriske parametrar. Dei fleste punkt i dette rommet er tomme; naturen har berre okkupert ein liten del av dei moglege formene. Spørsmålet som driv denne boka, er det same som dreiv Raup: kvifor er somme regionar busette og andre tomme?

Det empiriske formrommet er alltid ein projeksjon. Når me måler høgde, breidde og djupn på 2\,300 stolar og plottar dei i eit tredimensjonalt rom, har me ikkje fanga forma fullt ut. Me har projisert henne ned til tre dimensjonar. Den verkelege forma er rikare; ho rommar trekk som vinkelen på ryggen, krumlinja i setet, tjukkleiken på beina, og uendeleg mange andre eigenskapar me ikkje har målt. Inga endeleg mengd av målingar fangar alt.

Dette er ikkje ei svakheit ved metoden; det er ein eigenskap ved formverda. Projeksjonen vel kva relasjonar ho lyser opp og kva ho mørklegg. Det finst ingen nøytral projeksjon. Men forma eksisterer uavhengig av målinga; ho er der uavhengig av om nokon ser henne.

\section*{Klassen og algebraen}

\prop{1.31}{d}{Klassen avgrensar formrommet gjennom delt affordansestruktur: objekta i klassen tilbyr same konstellasjon av operasjonar til same konstellasjon av agentar. Klassen set grensene; seleksjonstrykket fordelar innhaldet.}

Kva avgrensar eit formrom? Kva avgjer at stolar og sofaer ikkje er same klasse? Det tradisjonelle svaret er: funksjonen. Stolen er til å sitje i; sofaen er til å liggje i. Men dette svaret er for grovt. Ein sofastol tilbyr begge delar. Ein krakk tilbyr sitting men ikkje ryggstøtte. Grensa mellom klassane er ikkje definert av ein enkelt funksjon, men av ein konstellasjon av affordansar: kva operasjonar objektet tilbyr til kva agentar.

Ein stol tilbyr sitting til ein voksen menneskekropp, stabling til ein kafétilsett, transport til eit fraktselskap, visuell karakter til eit sosialt rom. Desse affordansane saman avgrensar klassen. Eit objekt som deler desse affordansane, er ein stol; eit som ikkje gjer det, er noko anna.

Klassen set grensene for formrommet. Seleksjonstrykka; dei kreftene me skal definere i neste kapittel; fordeler innhaldet innanfor desse grensene. Klassen seier kvar rommet sluttar. Kreftene seier kvar i rommet formene samlar seg.

\prop{1.32}{d}{Morforommet er ein algebra lukka under union, differanse og transformasjon. Algebraen definerer kva operasjonar som er moglege.}

Formrommet er ikkje berre eit passivt sett av punkt. Det har algebraisk struktur: former kan kombinerast (union), trekkjast frå kvarandre (differanse), og transformerast (rotasjon, translasjon, skalering). Desse operasjonane er ikkje metaforar; dei er presise matematiske operasjonar definerte av George Stiny i hans arbeid med formalgebra og formgrammatikk.

\propcont{Lat S vere ei mengd av former i eit euklidsk rom E\textsuperscript{n}, utstyrt med den euklidske transformasjonsgruppa Trans(E\textsuperscript{n}):}

\begin{center}\small
\begin{tabular}{l@{\hspace{6mm}}l}
s\textsubscript{1} + s\textsubscript{2}   & union (sum) \\
s\textsubscript{1} \msym{−} s\textsubscript{2}   & differanse \\
s\textsubscript{1} \msym{·} s\textsubscript{2} := s\textsubscript{1} \msym{−} (s\textsubscript{1} \msym{−} s\textsubscript{2}) & snitt (avleidd) \\
\msym{τ}(a) \msym{≤} s, \quad \msym{τ} \msym{∈} Trans(E\textsuperscript{n}) & innleiring \\
\end{tabular}
\end{center}

\propcont{Ei \textbf{innleiring} av a i s er ein transformasjon som plasserer a i s, eventuelt rotert, flytta eller skalert. Innleiringa er ein eigenskap ved den noverande forma, uavhengig av korleis ho vart konstruert.}

Kvifor er algebraen viktig? Fordi han definerer kva som er moglege operasjonar i formrommet. Ein stolmakar som sveiser eit stålrøyr til ein eksisterande treramme, utfører ein algebraisk union. Ein designar som fjernar ornamentet frå ein historistisk stol, utfører ein algebraisk differanse. Ein fabrikk som skalerer ein prototype ned til barnestorleikar, utfører ein algebraisk transformasjon. Algebraen er ikkje ein teori om form; han er syntaksen til formrommet.

Innleiringa er den viktigaste operasjonen. Å oppdage at ei delform finst i ei anna form; at eit stolbein kan gjenkjennast i ein stolkonfigurasjon; er grunnlaget for formgrammatikken me skal definere i kapittel 5. Innleiringa opererer utelukkande på den noverande forma; ho bryr seg ikkje om korleis forma vart konstruert. Dette er ei avgjerande avklåring: grammatikken opererer i notid, uavhengig av produksjonshistorie.

\section*{Tre regionar}

\prop{1.4}{d}{Formrommet deler seg i tre regionar: busette, opne og forbodne.}

Ikkje alle posisjonar i formrommet er okkuperte. Dei fleste er tomme. Denne observasjonen er utgangspunktet for alt som følgjer. Dersom alle posisjonar var like sannsynlege; dersom det ikkje fanst krefter som favoriserte nokre konfigurasjonar over andre; ville formrommet vore uniformt busett. Det er det ikkje. Formene klumpar seg. Dei samlar seg i visse regionar og unngår andre. Denne klumpinga krev ei forklaring.

\prop{1.41}{o}{Dei busette regionane er det historiske arkivet. Dei fungerer som ankerpunkt; tyngda veks med okkupasjonstid. Uavhengig konvergens provar at optimumet er reelt.}

Dei busette regionane er dei posisjonane i formrommet der reelle konfigurasjonar finst. Det historiske arkivet; alle stolar som nokon gong er laga; er ein sky av punkt konsentrert i visse delar av rommet. I vårt korpus av 2\,300 europeiske stolar (1280--2024) frå Nasjonalmuseet og Victoria and Albert Museum, finn me to klare konsentrasjonar: éin kring (B \msym{≈} 50, H \msym{≈} 91) cm; den tradisjonelle høgrygga stolen; og éin kring (B \msym{≈} 50, H \msym{≈} 45) cm; den lågare modernistiske stolen. Mellom dei er det ein dal med få bebuarar.

Desse konsentrasjonane er ikkje tilfeldige. Dei er attraktorar i tilpassingslandskapet; eit omgrep me definerer i kapittel 3. At to uavhengige tradisjonar; den norske og den britiske; konvergerer mot same punkt i formrommet, er sterkare evidens for at attraktoren er reell enn at éin tradisjon åleine hadde gjort. Uavhengig konvergens er formverda sitt svar på eksperimentell replikasjon.

\prop{1.42}{o}{Dei opne regionane er det \textit{tilstøytande moglege}: tilgjengelege, men uaktualiserte. Dei er endelige; spente ut mellom det busette og det forbodne.}

Mellom dei busette regionane og dei forbodne ligg det opne territoriet: posisjonar som er fysisk moglege men som ingen har realisert. Stuart Kauffman kalla dette det «tilstøytande moglege» i sin teori om biologisk evolusjon. Det er innovasjonsradiusen; mengda av former som ligg éin operasjon unna dei formene som allereie eksisterer.

Det tilstøytande moglege er ikkje uendeleg. Det er avgrensa av dei busette regionane (som tiltrekkjer) og dei forbodne (som fråstøyter). Ein designar som arbeider med ein eksisterande stol, kan gjere ei lita endring; flytte ryggen framover, endre vinkelen på beina, byte materiale i setet; og dermed okkupere ein posisjon i det tilstøytande moglege. Men ho kan ikkje hoppe til ein vilkårleg posisjon i formrommet; avstandane er for store, og dei mellomliggjande posisjonane er tomme eller forbodne.

\prop{1.43}{d}{Dei forbodne regionane er utelukka av materielle, strukturelle eller energetiske vilkår. Forbodet følgjer vilkåra, ikkje regionen. Eit materialskifte kan opne ein forboden region.}

Ein stol med sete 3 meter over bakken er fysisk mogleg å byggje, men ingen gjer det. Tyngdekrafta, menneskekroppens mål og ergonomiske krav forbyr det. Ein stol i massivt gull er materialteknisk mogleg, men økonomiske seleksjonstrykk forbyr det. Ein stol med bein av papir er geometrisk mogleg, men materialet vetoar det.

Det vesentlege er at forbodet tilhøyrer vilkåra, ikkje regionen. Posisjonane i formrommet er i seg sjølve verken lovlege eller ulovlege; dei vert forbodne av dei materielle, strukturelle og energetiske restriksjonane som gjeld til kvar tid. Endrar du restriksjonane, endrar du grensa. Stål var ein forboden region for sitjemøbel i hundreår; damppressingsteknologien opna han. Plast var ein forboden region inntil sprøytestøyping gjorde han tilgjengeleg. Dei forbodne regionane er ikkje permanente; dei er dynamiske, og dynamikken deira er temaet for kapittel 4.

\section*{Tankerommet og kunnskapsrommet}

\prop{1.44}{d}{Dei tre regionane utgjer ein epistemisk partisjon: det \textbf{kjende} (K), det \textbf{tenkelege} (C\msym{∖}K), og det utenkjelege. Grensene er epistemiske, ikkje ontologiske.}

Dei tre regionane kan forståast som ein partisjon av kunnskap. Armand Hatchuel og Benoit Weil har formalisert dette i sin tanke--kunnskapsteori, som me her gjer til ein del av fundamentet.

C er tankerommet: alle former som er tenkelege men ikkje avgjorte. K er kunnskapsrommet: alle former som er kjende å verke eller svikte. Formgjeving er rørsla mellom dei.

\propcont{C er tankerommet: alle former som er tenkelege men ikkje avgjorte. K er kunnskapsrommet: alle former som er kjende å verke eller svikte. Formgjeving er rørsla mellom dei:}

\begin{center}\small
\begin{tabular}{l@{\hspace{4mm}}l@{\hspace{4mm}}l}
\textbf{Operator} & \textbf{Retning} & \textbf{Verknad} \\[2pt]
C\msym{→}C & partisjon & tanken vert spesifisert \\
C\msym{→}K & realisering & tanken vert avgjort \\
K\msym{→}K & deduksjon & ny kunnskap frå eksisterande \\
K\msym{→}C & konjunksjon & ny kunnskap opnar nye tankar \\
\end{tabular}
\end{center}

Tanken «ein stol utan bein» er i C; han er tenkeleg men ikkje avgjort. Verner Pantons eksperiment i 1960 flytta han til K; det viste seg å vere realiserbart. Denne realiseringa opna nye regionar i C: kva andre former kan eit sjølvberande plastskall ta? Det tilstøytande moglege ekspanderte.

Grensene mellom K, C\msym{∖}K og det utenkjelege er epistemiske, ikkje ontologiske. Dei seier ikkje kva som er mogleg i seg sjølv; dei seier kva me veit om kva som er mogleg. Eit nytt materiale; eit nytt verktøy; ein ny produksjonsprosess; kan flytte grensa mellom det tenkelege og det utenkjelege over natta. Heile formhistoria er historia om slike grenseforskyvingar.

\begin{overgang}
Kapittelet har definert formverda sine grunnlagande omgrep: objektet, konfigurasjonen, formrommet, algebraen, dei tre regionane og partisjonen mellom tanke og kunnskap. Neste kapittel stiller spørsmålet som driv alt: kvifor er somme posisjonar i formrommet meir sannsynlege enn andre?
\end{overgang}
